%!TEX ROOT=sample-thesis.tex
%!TEX PROGRAM=xelatex

\noindent
Kursus dalam talian terbuka besar-besaran (MOOCs) kini merupakan mod pembelajaran yang semakin lazim selaras dengan perkembangan platform pembelajaran seperti Coursera dan EdX. Platform pembelajaran dalam talian terkini menyediakan analitik pratetap yang membolehkan pengajar menganalisa data kursus mereka. Namun, kebolehgunaan visualisasi analitik pratetap untuk menganalisa data MOOC yang amat besar terlalu diterhadkan oleh mod analisis dan paparan tersedia. Literatur menunjukkan bahawa rekabentuk visualisasi MOOC tersedia hanya tertumpu kepada penganalisis MOOC, beserta perhatian terhad untuk membantu pengguna dominan iaitu pengajar. Kajian ini juga menemui sedikit rekabentuk visualisasi berlandaskan peranan pengajar dan kekerapan pemantauan kursus. Batasan ini dikaji dan menemukan jurang dalam rekabentuk visualisasi, terutamanya untuk analisis pemantauan sokongan pelajar. Demikian, kajian ini dilaksanakan untuk membantu analisis visual pengajar dalam pemantauan sokongan pelajar MOOC mengunakan kaedah berlandaskan permasalahan pengguna. Jurang kebolehgunaan rekabentuk dikenalpasti melalui analisa sistematik ruang rekabentuk visual analitik MOOC. Untuk mengatasi jurang tersebut, sebuah rekabentuk visual analitik berlandaskan pengguna MOOCSupport dicadangkan bagi membolehkan pemantauan sokongan pelajar secara visual. Lima pakar bidang telah ditemuramah untuk perincian masalah serta mengenalpasti keperluan rekabentuk. MOOCSupport ini dibina secara terperinci dengan mengambilkira kelebihan setiap visualisi secara sistematik bagi menghasilkan kebolehgunaan dan kemampuan belajar yang lebih baik. Kajian kes dan wawancara yang dilaksanakan bersama pakar bidang telah membuktikan kegunaan dan keberkesanan MOOCSupport dalam membantu analisis visual bagi pengajar. Sejurus itu, pereka visualisasi boleh memanfaatkan rekabentuk MOOCSupport beserta perincian proses rekabentuk yang dibentangkan bagi menyokong tenaga pengajar MOOC dalam pemantauan sokongan pelajar.
 
